\clearpage
\onecolumn
\section{Render Evidence Atlas}
\label{sec:supp-render-atlas}

This appendix front-loads the visual evidence that is otherwise scattered
across metric tables, GRScenes diagnostics, material-effect cases, and
navigation media. The goal is to give reviewers a compact way to inspect what
the paper means by a controlled visual perturbation: matched original-MDL and
noMDL renders, selected task-level target-view checks, and the selected
NVIDIA-baseline material comparison.

The page layout organizes existing panels from tracked render outputs and paper
figures. This keeps the supplement visually dense while preserving provenance.

\begin{center}
\includegraphics[width=\textwidth,height=0.58\textheight,keepaspectratio]{figures/fig_supplement_render_atlas_opener.png}
{\captionsetup{hypcap=false}
\captionof{figure}{Render evidence section opener. The compact panel uses
render crops to connect proxy pairs, GRScenes target views, material
diagnostics, and InternNav stills before the full-page atlas figures.}
\label{fig:supp-render-atlas-opener}}
\end{center}

\clearpage
\twocolumn

\begin{figure*}[p]
\centering
\includegraphics[width=\textwidth,height=0.86\textheight,keepaspectratio]{figures/fig_supplement_render_atlas.png}
\caption{Dense render atlas for the visual-preservation evidence chain. The
first four rows show proxy object pairs with both front context and contained
material-detail views under matched original-MDL and noMDL conditions. The
final two rows show paired GRScenes camera contexts. The atlas is qualitative
inspection material; metric tables and VLM/task tables provide the
quantitative context.}
\label{fig:supp-render-atlas}
\end{figure*}
\clearpage

\begin{figure*}[p]
\centering
\includegraphics[width=\textwidth,height=0.88\textheight,keepaspectratio]{figures/fig_supplement_render_scene_evidence_extended.png}
\caption{Extended render-scene evidence view used as an entry point from proxy
assets to scene-context evidence. The first two rows expose full proxy-object
context alongside contained material-detail views; the last two rows expose
matched GRScenes scene pairs. The same scope applies: these panels
show inspected qualitative preservation under selected cameras, while semantic
equivalence remains outside the measured result.}
\label{fig:supp-render-scene-wide}
\end{figure*}
\clearpage

\begin{figure*}[p]
\centering
\includegraphics[width=\textwidth,height=0.78\textheight,keepaspectratio]{figures/fig_grscene_qualitative.png}
\caption{GRScenes qualitative pair retained for visual continuity with the
main evidence chain. The paired views show how floor, wood, wall, and object
appearance change after noMDL conversion while preserving a usable scene frame
for downstream inspection.}
\label{fig:supp-grscene-qualitative-atlas}
\end{figure*}
\clearpage

\begin{figure*}[p]
\centering
\includegraphics[width=\textwidth,height=0.82\textheight,keepaspectratio]{figures/fig_supplement_vlm_clean_rerender_panel.png}
\caption{Post-repair GRScenes target-view render panel for the VLM evidence
chain. Each pair comes from a clean rerender smoke report with zero
original-MDL and zero converted noMDL MDL/KooPbr error signal; green boxes show
the target projections. The target close-up strip reuses the same
clean rerenders to make the boxes inspectable at page scale. No VLM
point marker is overlaid here, so the panel should be read as render-quality
support for the tabulated grounding evidence rather than as a rerun of the VLM
prediction study.}
\label{fig:supp-vlm-grounding-atlas}
\end{figure*}
\clearpage

\begin{figure*}[p]
\centering
\includegraphics[width=\textwidth,height=0.80\textheight,keepaspectratio]{figures/fig_material_effect_baseline_qualitative.png}
\caption{Selected material-effect comparison across original MDL,
ConvertAsset noMDL, and the NVIDIA baseline. Covered bins such as opacity or
transparency, emission, normal or bump, and displacement or height can be
visually inspected across all three columns. The figure supports selected
qualitative comparison; converter ranking would require a broader audit.}
\label{fig:supp-material-baseline-atlas}
\end{figure*}
\clearpage

\FloatBarrier
